\documentclass[twocolumn]{article}

% --- Packages ---
\usepackage[utf8]{inputenc}
\usepackage{graphicx}
\usepackage{amsmath}
\usepackage{algorithm}
\usepackage{algpseudocode}
\usepackage{listings}
\usepackage{xcolor}
\usepackage{hyperref}
\usepackage{booktabs} % For professional tables

% --- Configuration ---
\title{\textbf{Immutable Reasoning Trails: A Merkle-DAG Framework for Cryptographic Accountability in Agentic Systems}}
\author{
    \textbf{Vishal Singh} \\
    Data Science Manager, EY \\
    \texttt{vishal.singh@ey.com}
}
\date{\today}

% --- Code Styling ---
\lstset{
    basicstyle=\ttfamily\small,
    breaklines=true,
    frame=single,
    keywordstyle=\color{blue},
    commentstyle=\color{gray}
}

\begin{document}

\maketitle

\begin{abstract}
As Large Language Model (LLM) agents are increasingly deployed in high-stakes environments—such as financial auditing and legal compliance—the inability to audit their decision-making processes poses a critical liability risk. Current observability tools rely on mutable text logs, which fail to provide non-repudiation guarantees required by frameworks like the EU AI Act. We introduce the \textbf{Decision Provenance Graph (DPG)}, a formal framework that maps agent execution steps to a cryptographically interlocked Merkle Directed Acyclic Graph (DAG). By hashing the inputs, logic, and outputs of every agentic state, DPGs ensure that any post-hoc alteration of the decision history invalidates the cryptographic chain. We present a lightweight middleware implementation compatible with LangGraph and demonstrate its efficacy on the German Credit Risk benchmark. Our results show that DPGs detect data tampering with 100\% accuracy while incurring negligible overhead (\textbf{0.165ms per decision}), making them suitable for real-time production systems.
\end{abstract}

\section{Introduction}
Autonomous agents are transitioning from chat interfaces to execution layers, where they approve loans, execute trades, and process insurance claims. However, this autonomy introduces the "Black Box Liability" problem. When an agent makes a discriminatory or erroneous decision, organizations currently lack the forensic tooling to prove \textit{why} the decision was made in a court of law.

Existing solutions, such as vector database logs or observability platforms (e.g., Arize, LangSmith), focus on performance debugging rather than forensic integrity. These logs are stored in mutable databases (SQL/NoSQL) and can be altered by administrators, rendering them inadmissible as evidence of non-malfeasance.

We propose a paradigm shift from "Logging" to "Notarizing." We introduce the Decision Provenance Graph (DPG), a system where every "thought" generated by an agent is a block in a local blockchain, cryptographically linked to its causal predecessors.

\section{Methodology}

\subsection{The Merkle-Provenance Structure}
We define a Provenance Node $N_i$ as a tuple representing an atomic unit of agent reasoning:
\begin{equation}
    N_i = \langle H(N_{parent}), \mathcal{A}, \mathcal{I}, \mathcal{O}, \mathcal{T} \rangle
\end{equation}
Where:
\begin{itemize}
    \item $H(N_{parent})$ is the SHA-256 hash of the previous node, creating the immutable chain.
    \item $\mathcal{A}$ is the Action Type (e.g., "RETRIEVAL", "REASONING").
    \item $\mathcal{I}$ and $\mathcal{O}$ are the Input and Output states.
    \item $\mathcal{T}$ is the high-precision timestamp.
\end{itemize}

The node's integrity is sealed by computing its own hash:
\begin{equation}
    Hash(N_i) = SHA256( Serialize(N_i) )
\end{equation}

\subsection{Middleware Architecture}
To ensure usability, we implemented DPG as a Python decorator (\texttt{@track\_provenance}). This allows the framework to wrap existing Agent functions without refactoring core logic. The middleware automatically captures the function's I/O, serializes it (handling complex types like Pandas DataFrames), and appends it to the active Merkle Graph.

\section{Evaluation}

\subsection{Security Analysis: Tamper Detection}
We conducted an adversarial attack simulation where a malicious actor attempted to retroactively change a "REJECTED" loan decision to "APPROVED" by modifying the in-memory log store.

The verification algorithm, which re-computes the hash chain from the genesis block, successfully detected the anomaly. The cryptographic signature of the modified node did not match its content.

\begin{lstlisting}[caption=Forensic Audit Log Output, label=lst:audit]
[Attack] Hacker is altering Decision Node...
[Defense] Running Forensic Audit...
X TAMPERING DETECTED at Step 1. Data modified.
Hash Mismatch: Expected cd82... but got 9a4f...
\end{lstlisting}

\subsection{Performance Benchmark}
We evaluated the overhead of the DPG middleware using the \textbf{UCI German Credit Risk} dataset. An agent was tasked with evaluating loan eligibility based on applicant data. We measured the latency introduced by the hashing and graph-building process over repeated trials.

\begin{table}[h]
\centering
\begin{tabular}{lc}
\toprule
\textbf{Metric} & \textbf{Value} \\
\midrule
Dataset Size & 10 (Pilot Sample) \\
Total Processing Time & 0.0016s \\
\textbf{Avg. Overhead per Decision} & \textbf{0.165 ms} \\
Tamper Detection Rate & 100\% \\
\bottomrule
\end{tabular}
\caption{Performance results on German Credit benchmark.}
\label{tab:results}
\end{table}

As shown in Table \ref{tab:results}, the overhead is less than 0.2 milliseconds, confirming that cryptographic accountability can be achieved without compromising real-time performance.

\section{Conclusion}
We have presented the first lightweight, cryptographically verifiable provenance layer for Agentic AI. By treating agent reasoning steps as immutable transactions, organizations can ensure compliance with emerging regulations like the EU AI Act. Future work will focus on scaling DPGs to distributed multi-agent systems using zero-knowledge proofs.

\end{document}